\subsection{Классическая постановка задачи коммивояжера}

Пусть

\begin{equation}
    G = (V, E)
\end{equation}

--- полный неориентированный граф, где

\begin{itemize}
    \item $V = \{1, \dots, n\}$;
    \item $E = \{\{u, v\} | u, v \in V, u \neq v\}$.
\end{itemize}

Задана функция весов:
\begin{equation}
    d: V \times V \to \mathbb{R}_{\geq 0}
\end{equation}

такая, что выполняются условния:

\begin{itemize}
    \item симметрия: $\forall u, v \in V \quad d(u, v) = d(v, u)$;
    \item тождество нуля: $ \forall u \in v \quad d(u, u) = 0 $;
    \item неравенство треугольника: $ \forall u, v, w \in V \quad d(u, w) \leq d(u, v) + d(v, w) $.
\end{itemize}

Решением TSP называется перестановка графа $ \pi \in S_n $, минимизирующая значение

\begin{equation}
    cost(\pi) = \sum_{i=1}^{n-1} d(\pi(i), \pi(i+1)) + d(\pi(n), \pi(1))
\end{equation}

